A method described only in prose is hard to check and harder to reuse. Three
artifacts accompany this manuscript, and this section states what each
contains, what it can and cannot do, and where the boundary between the
deterministic engine and everything built on top of it lies. The boundary is
the design constraint that makes the rest possible: the engine has no
knowledge of the web layer, and every array crossing into it is the same array
the experiments in Section~\ref{sec:results} were run on.
Figure~\ref{fig:artifact-stack} shows the arrangement.

\begin{figure}[htbp]
\centering
\begin{tikzpicture}[node distance=5.5mm and 12mm]
  \node[fcstart, minimum width=68mm] (engine)
    {\textbf{Deterministic engine}\\[1pt]
     \scriptsize CST algebra, constraint rows, pre-solve, extended Newton
     solver, diagnostics\\[1pt]
     \scriptsize no input/output, no web dependency; the code the experiments ran};
  \node[fcbox, below=of engine, minimum width=68mm] (api)
    {\textbf{HTTP service}\\[1pt]
     \scriptsize one adapter module is the only caller of the engine;
     routers translate requests and responses only};
  \node[fcout, below=9mm of api, xshift=-17mm, minimum width=32mm] (web)
    {\textbf{Interactive application}\\[1pt]
     \scriptsize analyse, inverse design,\\[-1pt]
     \scriptsize flow field, gallery};
  \node[fcout, below=9mm of api, xshift=17mm, minimum width=32mm] (site)
    {\textbf{Project site}\\[1pt]
     \scriptsize narrative walkthrough,\\[-1pt]
     \scriptsize figures and progress record};
  \node[fcbox, below=22mm of api, minimum width=68mm, fill=black!3] (repo)
    {\textbf{Repository}: source, configurations, run records, this manuscript};
  \draw[fcarrow] (engine) -- (api);
  \draw[fcarrow] (api.south) -- ++(0,-3mm) -| (web.north);
  \draw[fcarrow] (api.south) -- ++(0,-3mm) -| (site.north);
  \draw[fcarrow, dashed] (repo.west) -- ++(-6mm,0) |- (engine.west);
\end{tikzpicture}
\caption{The released artifacts. The engine is the code the experiments ran;
everything above it is an adapter, and nothing above it can change a result.}
\label{fig:artifact-stack}
\end{figure}

\subsection{The interactive application}

A deployed web application~\citep{cinsapp} exposes the method to a reader who
does not want to install anything. It is a Python service wrapping the engine,
with a browser front end, and it offers the four views shown in
Figure~\ref{fig:app-views}. \emph{Analyse} runs a direct
coupled solve for a section chosen from the same 123-file coordinate corpus
the panel of Section~\ref{sec:h1-uiuc-panel} was drawn from, or from a NACA
code, and reports the pressure distribution, the boundary-layer state, and the
CST-derived engineering parameters, being leading-edge radius, trailing-edge
wedge half-angles, maximum thickness and camber with their locations, and
inscribed area, all evaluated in closed form from the coefficients by the same
Beta-function machinery as Eq.~\eqref{eq:area-beta} rather than by quadrature.
\emph{Inverse design} runs the monolithic solve itself, iteration by
iteration, in either of two modes: a self-consistency solve against a target
generated from a section's own fit, which is the experiment of
Section~\ref{sec:falsifiable}, or a solve against a target the user supplies
as a pressure or edge-velocity curve. \emph{Flow field} renders the inviscid
velocity and pressure field around the section. \emph{Gallery} replays the
archived results and figures this manuscript reports.

Two behaviours of the user-supplied-target mode are worth recording here,
because they are the product-side evidence for
Section~\ref{sec:le-identifiability}'s finding. First, the realisability
verdict is always reported and never blocks: a target outside the
CST-representable manifold produces a warning carrying the
inviscid-consistent metric of Section~\ref{sec:presolve}, not an error, which
is the reporting rule of that section made into interface behaviour. Second,
the same station-restriction rule the panel runs use is applied here, and it
was in fact found here first. On a user-defined viscous target the
unrestricted candidate set converged in 14 iterations to a surface $0.86$
millichords from the target, while restricting candidates to stations aft of
the prescribed nose converged in 8 to $0.017$ millichords, a fiftyfold
improvement in recovered shape at an almost unchanged residual. Recovery on
such a target is reported and asserted on the \emph{shape} rather than on the
coefficients, because the CST basis is ill-conditioned in $A$
(Section~\ref{sec:nsweep}) and coefficient vectors differing by
$1.8\x10^{-2}$ describe surfaces differing by less than half a millichord;
against a target that only ever constrained pressure, the coefficient norm is
not a meaningful accuracy statement.

\begin{figure}[htbp]
\centering
\includegraphics[width=\textwidth]{paper/fig_app_views.png}
\caption{The four views of the deployed application, captured from the running
service: (a)~analyse; (b)~inverse design; (c)~flow field; (d)~gallery.}
\label{fig:app-views}
\end{figure}

Figure~\ref{fig:app-views} is worth reading against the formulation rather than
as a product tour. Panel~(a) shows the engineering parameters of
Eq.~\eqref{eq:endpoint} and Eq.~\eqref{eq:area-beta} evaluated on a solved
section, which is the closed-form geometric-functional machinery of
Section~\ref{sec:constraints} exposed directly rather than inferred from
coordinates. Panel~(b) is the inverse problem as the user states it, a baseline
geometry, a target pressure curve, and an optional constraint row, which is
precisely the three-block decomposition of Eq.~\eqref{eq:jacobian-blocks}
presented as a form; the constraint panel is also where the target-consistent
right-hand-side rule of Appendix~\ref{app:area} is enforced, since a
leading-edge row added there is built from the pre-solved guess rather than
from an idealised value. Panel~(c) shows the inviscid field beside the coupled
boundary layer, which is the split Section~\ref{sec:transition} depends on.
Panel~(d) replays archived runs, so the evidence behind
Section~\ref{sec:results} can be inspected without a solve.

Two limitations of the deployment are stated so that a reader who tries it is
not misled. The service runs one inverse solve at a time by construction,
because the forced-transition closure of Section~\ref{sec:transition} replaces
module-level functions in the vendored solver and is therefore process-global
rather than instance-scoped; a second submission waits rather than running
concurrently. And the rendered flow field is inviscid, since the
displacement-thickness source contribution is not yet added to it, while the
boundary-layer quantities displayed beside it come from the fully coupled
solve.

\subsection{The project site and the repository}

A project site~\citep{cinssite} presents the method as a narrative walkthrough, from the
problem statement through the linearity argument, the square system, the
falsifiable test, the identifiability finding and the conditioning cliff, using
the same figures as this manuscript. It exists because the argument of
Section~\ref{sec:cst-linearity} is visual before it is algebraic, and a reader
deciding whether to spend time on the formulation is better served by the
figures than by the equations.

The repository carries the engine, the configuration files that define every
run reported here, the archived run records those runs produced, the figure
scripts, the test suite, and the source of this manuscript.
Appendix~\ref{app:repro} states the reproduction procedure and the software
environment; Appendix~\ref{app:prov} maps each figure and table in this
manuscript to the record it was produced from. The vendored flow solver is
included unmodified under its own licence and is never edited; all access to
it passes through a single adapter module, which is what made the
forced-transition closure of Section~\ref{sec:transition} implementable
without a fork.
